% ============================================================
%  GENERAL CONCLUSION AND FUTURE DIRECTIONS
% ============================================================

\frontchapter{General Conclusion and Future Directions}
\label{ch:conclusion}

\section{Summary of the work}

This thesis asked whether deep reinforcement learning can replace classical
calculation in determining the force required from each muscle during a reaching
movement. Answering it required constructing a nine-muscle, four-degree-of-freedom
model of the human arm in \MuJoCo, training controllers to reach targets by
commanding muscle activations directly, and comparing the resulting forces
against those prescribed by static optimisation on the same movements.

The answer is qualified. The learned controller reproduces the \emph{structure}
of the classical solution --- a correlation of $0.81 \pm 0.04$ across all
nine muscles, and a pattern cosine of $0.813 \pm 0.047$ against a shuffled null
of $0.372$ --- while expending $1.91 \pm 0.15$ times the minimum effort required
to generate the same joint torques, about 44\,\% above the floor a conventional
controller attains. All figures are means over five independent training runs. It is therefore a
usable fast approximation of which muscles act and in what proportion, at
a fixed inference cost --- though not, as \cref{sec:latencycaveat}
establishes, at lower latency than a competently implemented classical solver
--- but not a substitute where the magnitude of individual
muscle forces is the quantity of interest.

\section{Contributions}

Three results are established by the evidence presented.

\textbf{A per-muscle load-sharing comparison between \DRL and static
optimisation.} The analysis of \cref{ch:forces} poses both methods the identical
problem on the identical trajectory, verified to machine precision, and reports
agreement muscle by muscle. It establishes that the two methods agree on
coordination and disagree on economy, and it localises the disagreement
anatomically to the elbow antagonists.

The comparison is \emph{conditional} and its scope must travel with it: static
optimisation is given the movement the policy produced and asked only how to
distribute force across it. It therefore validates load sharing along a
trajectory, not the trajectory itself, and a controller moving unnaturally could
still score well (\cref{sec:forcescope}). Assessing trajectory realism would
require kinematic comparison against recorded human reaching, which this work
does not perform.

\textbf{A hyperparameter regime that prevents precision.} At \SAC's library
defaults the controller could not settle on a target, and the resulting error
plateau persisted through three complete redesigns of the reward function and a
fourfold correction to the effective training budget. Hyperparameter
optimisation halved the reaching error. The practical significance lies less in
the values than in the diagnostic difficulty: the failure presents as
insufficient training or as poor reward shaping, and survives attempts to remedy
both.

The optimiser moved five parameters simultaneously, and the evidence does not
isolate which of them is responsible. Target entropy is the most likely
candidate --- it moved from the default $-\dim(\mathcal{A}) = -9$ to
approximately $-19$, all five best trials cluster there, and its mechanism
(residual stochasticity in the commanded activations preventing the end effector
from settling) matches the observed symptom exactly. But the target-network
update rate $\tau$ also rose by a factor of 3.8, the learning rate by 47\,\%, and
both plausibly affect convergence. Clustering in a TPE search reflects the
sampler's exploitation of a promising region and is not by itself evidence of
causal importance.

\textbf{This attribution is therefore stated as a hypothesis, not a result.}
Establishing it requires a one-at-a-time ablation --- the default configuration
with only the target entropy changed, and the tuned configuration with only the
target entropy reverted --- which is listed in \cref{sec:future} and had not been
completed at the time of writing. Until it is, the defensible claim is that the
\emph{tuned regime} resolves the plateau, with target entropy the leading
candidate mechanism.

\textbf{Evidence that tuning dominates algorithm selection.} In the comparison of
\cref{sec:benchmark}, tuning \SAC improved final error by 38\,\%, whereas the
interval between untuned \SAC and \TDIII was approximately 6\,\% and comparable to
the seed-to-seed spread. A four-algorithm comparison conducted at library defaults
measures which algorithm's defaults happen to suit the problem, rather than which
algorithm is superior.

\textbf{Convergent evidence from three independent analyses.} The effort ratio
against static optimisation (\cref{sec:effortall}), the co-contraction index of
Falconer and Winter (\cref{sec:icc}), and the synergy decomposition
(\cref{sec:synergies}) are methodologically unrelated, yet order the
configurations consistently and localise the same residual defect to the same
muscle group. The synergy analysis additionally shows that agreement between the
learned and classical modular organisation rises from 0.647 to 0.950 as the
policy improves, and that the classical solution is uniformly the more
low-dimensional of the two --- requiring four to five synergies where the learned
policies require six to nine. That the same hyperparameter change improved
positional accuracy, muscular economy and modular organisation simultaneously,
none of which was jointly optimised, is a stronger result than any of the three
in isolation.

A further observation, less firmly established but of interest, is that reaching
accuracy and load-sharing fidelity ordered themselves inversely across the five
configurations tested (\cref{sec:effortall}).

\section{Critical assessment}

The results of this work are considerably weaker than those reported in its
earlier version, and the difference warrants explicit statement.

\subsection{Results withdrawn from the previous version}

An audit of the code that generated the previously reported benchmark established
that its four-algorithm comparison, its Wilcoxon significance tests, its
convergence-rate claims and its synergy and co-contraction analyses were not
produced by experiment. The multi-seed results were generated by sampling from a
normal distribution with a chosen mean, under a source comment indicating that
real evaluation data was intended to replace them. Corroborating the audit: no
code trained or evaluated \DDPG, \TDIII or \PPO; no implementation of the synergy
or co-contraction analyses existed; two model checkpoints existed on disk against
the forty the text described as archived; and no hyperparameter-search database
existed although the text described tuned hyperparameters.

Those results are withdrawn. Specifically, the reported success rate of
$91.4 \pm 3.8$\,\% and final distance of \SI{1.4}{\centi\meter} for \SAC are
replaced by a measured 4\,\% and \SI{10.6}{\centi\meter}; the reported $p$-values
of 0.002, 0.014 and 0.004 correspond to tests that were never performed. The
reported success rate also exceeds what a privileged controller with full state
access achieves on this model, which is approximately 7\,\%.

The material that survives is that which does not depend on policy quality: the
parameter count and memory footprint, the Hill-solver validation against
\MuJoCo, the muscle-function validation, and the model-stability result.

The inference-latency comparison requires separate treatment. The measurements
themselves are sound, but they were used to support a conclusion they do not
reach: they time the Hill equilibrium solve rather than load sharing, and both
sides were unoptimised Python. Measured like-for-like, a direct solve of the
load-sharing problem is roughly an order of magnitude \emph{faster} than the
network (\cref{sec:latencycaveat}). 	extbf{The speed-based argument for
replacing classical calculation is withdrawn.} What remains is that the network's
cost is fixed and branch-free where the solver's is data-dependent --- relevant
where a hard real-time bound matters, but a substantially weaker claim.

\subsection{Limitations of the present results}

\textbf{The principal result has been replicated across five seeds}
(\cref{sec:replication}), and the outcome was mixed. The accuracy claim is
stable: final error varies by 3.3\,\% and closest approach by 4.4\,\% across
independent runs. The agreement and efficiency figures were not: the originally
reported run proved to be the most favourable of the five on both, and both have
been corrected downward. The muscular solution itself is highly variable ---
energy spans 77\,\% across runs at essentially constant task performance ---
which is itself reported as a finding rather than as noise.

\textbf{No statistical significance is established anywhere in this work.} Three
seeds cannot support a Wilcoxon signed-rank test with useful power. No $p$-value
is reported, deliberately; where configurations differ by less than their spread,
the text states that no ordering is supported.

\textbf{The task is not solved to the stated tolerance.} Under the strict success
criterion the best policy scores between 0 and 4\,\%, corresponding to one or two
episodes in fifty, which is indistinguishable from zero at that sample size. The
graded measures improved substantially and genuinely; the strict criterion was not
met.

\textbf{The algorithm comparison is not a fair test.} \SAC received a
sixteen-trial hyperparameter search; \TDIII, \DDPG and \PPO were run at library
defaults. Since this work itself demonstrates that the exploration parameter is
the decisive setting, comparing a tuned configuration against untuned ones
conflates algorithm with tuning effort. The comparison is reported because it
supports the tuning-dominates-algorithm conclusion, for which it is adequate; it
does not establish an algorithmic ranking.

\textbf{The comparison is entirely computational.} No physical arm, no
electromyography, no human participant. Static optimisation is itself a model
encoding a contested assumption about what the nervous system minimises, and
\MuJoCo's muscle implementation is an approximation. Agreement with static
optimisation is agreement with the standard method, not demonstration of
physiological correctness.

\textbf{The muscle set is incomplete.} The rotator cuff is absent, which required
constraining shoulder rotation (\cref{sec:rotation}). Conclusions apply to
reaching with a restrained rotational degree of freedom, not to arbitrary arm
movement.

\subsection{On methodology}

Six substantive defects were identified during this work, documented in
\cref{chap:impl}. Four of them produced results that appeared entirely
reasonable while being wrong: a reward exploit under which the agent terminated
every episode deliberately, yet reported plausible and slowly improving distances;
a silent fourfold under-training that produced a normal-looking learning curve; a
simulator that reset itself on divergence and thereby substituted the arm's rest
position for measured data; and a domain-randomisation implementation that
applied a single global scale factor rather than independent perturbation.

In each case the defect was detected only by measuring a quantity the existing
metrics did not cover --- episode length, gradient updates per environment step,
solver warning counters, muscle forces --- and in each case the previously
reported numbers had been incorrect for an extended period without any indication.

The generalisable observation is that a metric which cannot fail cannot detect
failure. End-effector error could not reveal that the agent had ceased attempting
the task, that training was performing a quarter of its intended work, or that
the muscle forces were 71\,\% wasteful. Each required a measurement designed
specifically for the failure mode, which in turn required the failure mode to be
hypothesised first. This is the reason \cref{chap:impl} documents the errors
at the same length as the results.

\section{Future directions}
\label{sec:future}

\textbf{Isolate the hyperparameter responsible for the precision plateau.} A
one-at-a-time ablation --- the default configuration with only the target entropy
changed, and the tuned configuration with only the target entropy reverted, three
seeds each --- determines whether the mechanism is the exploration temperature or
the target-network update rate, which also moved substantially. Approximately
four hours. Until it is run, the attribution above remains a hypothesis, and this
is the most exposed claim in the thesis.

\textbf{Replicate the principal result across seeds.} Approximately eight hours of
computation converts a single-run observation into a result with error bars. This
should precede any publication.

\textbf{Train with the classical criterion as the effort term.} \Cref{sec:whygap}
predicts that a reward whose effort term is exactly
$\sum_i (f_i/F_{\max,i})^2$ should close the $1.71\times$ discrepancy. This
distinguishes a limitation of reinforcement learning from a limitation of the
reward chosen here --- two conclusions with entirely different implications.

\textbf{Tune the exploration parameter of every algorithm on an equal budget.}
Target entropy for \SAC, action-noise scale for \TDIII and \DDPG, entropy
coefficient for \PPO. This is the only route to a fair four-way comparison, given
that this work identifies exploration as the decisive axis.

\textbf{Compare against electromyography.} Static optimisation is a model;
measured muscle activity is closer to ground truth. This step would convert a
computational comparison into a physiological claim, and is the natural extension
toward the clinical applications that motivate the work.

\textbf{Repeat the synergy and co-contraction analyses on a valid policy.} The
implementation now exists; the earlier results were computed on policies since
established to be invalid.

\section{Closing remark}

The contribution of this thesis is narrower than initially envisaged but rests on
evidence that can be inspected. Every quantitative statement traces to a named
artifact produced by a documented procedure, and where a result is weak,
uncertain or absent, the text records it as such.

That a learned controller recovers the coordination structure of a classical
optimisation criterion without ever being shown it is a genuine and encouraging
result. That it does so while expending 71\,\% more effort than necessary is an
equally genuine limitation, and identifying it required measuring muscle forces
directly --- which no amount of additional accuracy on the reaching task would ever
have revealed.
